%
En leketøyskanon har en fjær med fjærkonstant $k = \SI{500}{\newton\per\metre}$. Fjæra presses sammen $x = \SI{8.0}{\centi\metre}$, og kanonen skyter en kule med masse $m = \SI{50}{\gram}$ rett opp. Se bort fra luftmotstand og fra den lille høydeendringen mens fjæra strekker seg ut.
\begin{enumerate}[a)]
    \item Hvor mye energi er lagret i fjæra før skuddet?
    \item Hvor høyt over utskytingspunktet kommer kula?
    \item Hvor høyt kommer kula dersom fjæra bare presses sammen halvparten så mye?
\end{enumerate}
